% ============================================================
% CHAPITRE 14 — Introduction à Adobe Illustrator
% ============================================================

\chapter{Introduction à Adobe Illustrator}

\noindent
\textit{Nous voici au début de la troisième partie de ce cours, consacrée à la maîtrise d'Adobe Illustrator. Dans les chapitres précédents, nous avons exploré en profondeur Photoshop : retouche, photomontage, texte, effets, création de supports publicitaires. Photoshop est un outil indispensable pour tout designer graphique, mais il a ses limites : il travaille sur des images matricielles, composées de pixels, qui perdent de la qualité à l'agrandissement. Pour créer des logos, des illustrations vectorielles, des mises en page précises, il faut un autre outil : Adobe Illustrator.}

\vspace{1em}

\noindent
Adobe Illustrator est le \terme{standard professionnel} de la création vectorielle. Il est utilisé par les designers graphiques, les illustrateurs, les typographes, les motion designers, les designers d'interface. Il permet de créer des logos, des illustrations, des pictogrammes, des icônes, des mises en page, des infographies, des typographies. Ses images vectorielles sont composées de formes géométriques (points, lignes, courbes), et non de pixels : elles peuvent être agrandies à l'infini sans perte de qualité.

\vspace{1em}

\noindent
Illustrator et Photoshop sont complémentaires. Photoshop excelle dans la retouche et la manipulation des photographies ; Illustrator excelle dans la création de formes nettes, précises, redimensionnables. Un logo se crée dans Illustrator, une affiche photographique dans Photoshop. Un designer professionnel maîtrise les deux, et sait choisir le bon outil pour le bon usage.

\vspace{1em}

\noindent
Ce chapitre a pour objectif de te faire découvrir Adobe Illustrator de manière progressive et structurée. Nous commencerons par une présentation générale du logiciel et de ses domaines d'utilisation. Nous étudierons ensuite les notions fondamentales d'image vectorielle, et la différence avec les images matricielles. Nous verrons l'interface du logiciel, les plans de travail, les outils essentiels, la création et l'enregistrement d'un document, les raccourcis clavier, et enfin les formats AI, SVG, EPS et PDF. À la fin de ce chapitre, tu disposeras des bases nécessaires pour utiliser Illustrator de manière efficace et professionnelle.

% ============================================================
\section{Présentation d'Adobe Illustrator}
% ============================================================

\subsection{Qu'est-ce qu'Adobe Illustrator ?}

\begin{quote}
\textbf{Adobe Illustrator est un logiciel de création d'images vectorielles, développé par Adobe Inc. Il permet de créer des logos, des illustrations, des pictogrammes, des mises en page, des typographies, des infographies, avec une précision et une qualité professionnelles.}
\end{quote}

Illustrator est un logiciel de \terme{PAO} (Publication Assistée par Ordinateur), plus précisément un logiciel de \terme{création vectorielle}. Il travaille sur des images composées de formes géométriques (points, lignes, courbes), contrairement à Adobe Photoshop qui travaille sur des images composées de pixels.

\subsection{Un peu d'histoire}

Illustrator a été créé en \terme{1987} par Adobe. La première version, \terme{Illustrator 1.0}, était destinée au Macintosh. Le logiciel a été conçu pour créer des illustrations vectorielles de qualité professionnelle, en s'appuyant sur le langage \terme{PostScript}, développé par Adobe.

Depuis, Illustrator a connu de nombreuses versions, chacune apportant de nouvelles fonctionnalités :

\begin{itemize}
    \item \terme{1987} : Illustrator 1.0, première version.
    \item \terme{1989} : Illustrator 3.0, introduction du texte.
    \item \terme{1993} : Illustrator 5.0, introduction des calques.
    \item \terme{1998} : Illustrator 8.0, introduction des dégradés de formes.
    \item \terme{2003} : Illustrator CS (Creative Suite), intégration à la suite Adobe.
    \item \terme{2013} : Illustrator CC (Creative Cloud), passage à l'abonnement.
    \item \terme{Aujourd'hui} : versions régulières avec nouvelles fonctionnalités (IA, etc.).
\end{itemize}

Illustrator est aujourd'hui utilisé par des millions de professionnels et d'amateurs dans le monde entier.

\subsection{Les caractéristiques principales}

\subsubsection{Logiciel vectoriel}

Illustrator travaille sur des \terme{images vectorielles}, composées de formes géométriques. Cela permet une qualité parfaite à toutes les tailles, sans pixellisation (voir section 3).

\subsubsection{Logiciel de création}

Illustrator excelle dans la \terme{création} : logos, illustrations, pictogrammes, typographies, infographies. Il est l'outil de référence pour les identités visuelles.

\subsubsection{Logiciel de précision}

Illustrator permet une \terme{précision} extrême : on peut positionner un élément au pixel près, définir des dimensions exactes, aligner avec rigueur.

\subsubsection{Logiciel professionnel}

Illustrator est un \terme{standard professionnel}, utilisé dans les agences, les studios, les entreprises. Sa maîtrise est souvent exigée dans les offres d'emploi.

\subsection{Illustrator dans la suite Adobe}

Illustrator fait partie de la \terme{suite Adobe Creative Cloud}, qui comprend également :

\begin{itemize}
    \item \terme{Adobe Photoshop} : retouche photographique, images matricielles.
    \item \terme{Adobe InDesign} : mise en page éditoriale.
    \item \terme{Adobe After Effects} : motion design.
    \item \terme{Adobe Premiere Pro} : montage vidéo.
    \item \terme{Adobe XD / Figma} : design d'interface.
    \item \terme{Adobe Lightroom} : retouche photographique.
\end{itemize}

Ces logiciels sont complémentaires : Illustrator travaille sur les images vectorielles, Photoshop sur les images matricielles, InDesign sur la mise en page. Un designer professionnel maîtrise souvent plusieurs de ces outils.

\subsection{Les versions d'Illustrator}

Illustrator existe en plusieurs versions :

\begin{itemize}
    \item \terme{Illustrator CC} : version complète, par abonnement (Creative Cloud).
    \item \terme{Illustrator iPad} : version pour tablette.
    \item \terme{Illustrator Draw} : version simplifiée pour mobile.
\end{itemize}

Dans ce cours, nous utiliserons la version \terme{Illustrator CC}, qui est la version professionnelle.

\subsection*{Conclusion de la section}

Adobe Illustrator est un logiciel de création d'images vectorielles, développé par Adobe. Créé en 1987, il est devenu un standard professionnel. Il permet de créer des logos, des illustrations, des pictogrammes, des typographies, des infographies. Il fait partie de la suite Adobe Creative Cloud, aux côtés de Photoshop, InDesign, After Effects, etc. La version utilisée dans ce cours est Illustrator CC.

Maintenant que nous avons présenté Illustrator, nous allons, dans la section suivante, étudier ses \terme{domaines d'utilisation}.

% ============================================================
\section{Les domaines d'utilisation d'Illustrator}
% ============================================================

\subsection{Pourquoi connaître les domaines d'utilisation ?}

Illustrator est un logiciel très polyvalent, utilisé dans de nombreux domaines. Connaître ces domaines permet de :

\begin{itemize}
    \item \terme{Comprendre} l'étendue des possibilités du logiciel.
    \item \terme{Choisir} les fonctionnalités adaptées à son usage.
    \item \terme{Se positionner} professionnellement.
    \item \terme{Développer} ses compétences dans les domaines pertinents.
\end{itemize}

\subsection{Les principaux domaines d'utilisation}

\subsubsection{La création de logos}

C'est le domaine historique d'Illustrator. Il permet de créer des \terme{logos} professionnels :

\begin{itemize}
    \item \terme{Logos typographiques} : nom écrit dans une typographie caractéristique.
    \item \terme{Logos symboliques} : symbole graphique.
    \item \terme{Logos mixtes} : combinaison de typographie et de symbole.
    \item \terme{Logos abstraits} : forme géométrique.
\end{itemize}

Illustrator est l'outil de référence pour les logos, car il permet de créer des formes nettes, précises, redimensionnables à l'infini.

\subsubsection{L'identité visuelle}

Illustrator permet de créer l'ensemble des éléments d'une \terme{identité visuelle} :

\begin{itemize}
    \item Logo et variantes.
    \item Palette de couleurs.
    \item Typographies.
    \item Éléments graphiques secondaires (formes, motifs, icônes).
    \item Charte graphique.
    \item Déclinaisons (papeterie, signalétique, web, packaging).
\end{itemize}

\subsubsection{L'illustration vectorielle}

Illustrator permet de créer des \terme{illustrations vectorielles} :

\begin{itemize}
    \item \terme{Illustrations éditoriales} : pour la presse, l'édition.
    \item \terme{Illustrations publicitaires} : pour les campagnes.
    \item \terme{Illustrations techniques} : schémas, plans, modes d'emploi.
    \item \terme{Illustrations jeunesse} : albums, livres pour enfants.
    \item \terme{Illustrations web} : icônes, pictogrammes, infographies.
\end{itemize}

\subsubsection{Les pictogrammes et icônes}

Illustrator est utilisé pour créer des \terme{pictogrammes} et des \terme{icônes} :

\begin{itemize}
    \item \terme{Signalétique} : panneaux, transports, espaces publics.
    \item \terme{Interfaces} : icônes d'applications, de sites.
    \item \terme{Infographies} : pictogrammes de données.
    \item \terme{Jeux vidéo} : icônes d'interface.
\end{itemize}

\subsubsection{La typographie}

Illustrator permet de créer et de manipuler de la \terme{typographie} :

\begin{itemize}
    \item \terme{Compositions typographiques} : titres, slogans, accroches.
    \item \terme{Lettrage} : personnalisation de lettres.
    \item \terme{Typographie sur tracé} : texte sur une courbe.
    \item \terme{Déformation de texte} : effets typographiques.
    \item \terme{Création de polices} : développement de caractères.
\end{itemize}

\subsubsection{Les infographies}

Illustrator est utilisé pour créer des \terme{infographies} :

\begin{itemize}
    \item \terme{Graphiques} : histogrammes, camemberts, courbes.
    \item \terme{Schémas} : processus, organigrammes.
    \item \terme{Cartes} : cartographies, plans.
    \item \terme{Visualisations} : données, statistiques.
\end{itemize}

\subsubsection{La mise en page}

Illustrator permet de créer des \terme{mises en page} simples :

\begin{itemize}
    \item \terme{Affiches} : compositions graphiques.
    \item \terme{Dépliants} : brochures, flyers.
    \item \terme{Couvertures} : livres, magazines, albums.
    \item \terme{Packaging} : étiquettes, emballages.
\end{itemize}

Attention : pour la mise en page éditoriale complexe (livres, magazines de plusieurs pages), on utilise plutôt InDesign.

\subsubsection{Le web et le digital}

Illustrator permet de créer des éléments pour le \terme{web} :

\begin{itemize}
    \item \terme{Icônes} : pour les interfaces.
    \item \terme{Logos} : pour les sites web.
    \item \terme{Illustrations} : pour les pages web.
    \item \terme{Infographies} : pour les articles.
    \item \terme{Visuels} : pour les réseaux sociaux.
\end{itemize}

Les fichiers SVG, exportés depuis Illustrator, sont parfaits pour le web : légers, redimensionnables, modifiables.

\subsubsection{Le motion design}

Illustrator est utilisé pour créer des éléments destinés au \terme{motion design} :

\begin{itemize}
    \item \terme{Illustrations} : pour les animations.
    \item \terme{Logos} : pour les génériques.
    \item \terme{Icônes} : pour les interfaces animées.
    \item \terme{Personnages} : pour les dessins animés.
\end{itemize}

Les fichiers Illustrator peuvent être importés dans After Effects pour être animés.

\subsubsection{Le packaging}

Illustrator est utilisé pour créer des \terme{packagings} :

\begin{itemize}
    \item \terme{Étiquettes} : pour les produits.
    \item \terme{Emballages} : boîtes, flacons, sachets.
    \item \terme{Gabarits} : découpes, pliages.
    \item \terme{Maquettes} : simulations 3D.
\end{itemize}

\subsection{Tableau récapitulatif des domaines}

\begin{table}[h!]
\centering
\begin{tabular}{@{}p{4cm}p{10cm}@{}}
\toprule
\textbf{Domaine} & \textbf{Utilisation principale} \\
\midrule
Logos & Création de logos et variantes \\
Identité visuelle & Systèmes graphiques complets \\
Illustration & Illustrations éditoriales, publicitaires, techniques \\
Pictogrammes & Signalétique, interfaces, infographies \\
Typographie & Compositions, lettrage, polices \\
Infographies & Graphiques, schémas, cartes \\
Mise en page & Affiches, dépliants, couvertures, packaging \\
Web & Icônes, logos, illustrations, SVG \\
Motion design & Éléments pour l'animation \\
Packaging & Étiquettes, emballages, gabarits \\
\bottomrule
\end{tabular}
\caption{Les principaux domaines d'utilisation d'Illustrator.}
\end{table}

\subsection{Les limites d'Illustrator}

Bien que très polyvalent, Illustrator a des limites :

\begin{itemize}
    \item \terme{Retouche photographique} : Illustrator n'est pas adapté à la retouche de photos (Photoshop est plus adapté).
    \item \terme{Photomontage} : Illustrator n'est pas adapté au photomontage (Photoshop est plus adapté).
    \item \terme{Mise en page éditoriale} : Illustrator n'est pas adapté aux documents longs (InDesign est plus adapté).
    \item \terme{Motion design} : Illustrator n'est pas adapté à l'animation (After Effects est plus adapté).
    \item \terme{Design d'interface} : Illustrator n'est pas le meilleur outil pour les interfaces (Figma, Sketch sont plus adaptés).
\end{itemize}

Un designer professionnel utilise le bon outil pour le bon usage.

\subsection*{Conclusion de la section}

Illustrator est utilisé dans de nombreux domaines : création de logos, identité visuelle, illustration vectorielle, pictogrammes, typographie, infographies, mise en page, web, motion design, packaging. Chaque domaine mobilise des compétences spécifiques. Illustrator a également des limites : il n'est pas adapté à la retouche photo, au photomontage, à la mise en page éditoriale, au motion design, au design d'interface. Un designer professionnel utilise le bon outil pour le bon usage.

Maintenant que nous connaissons les domaines d'utilisation d'Illustrator, nous allons, dans la section suivante, étudier les \terme{images vectorielles}.

% ============================================================
\section{Comprendre les images vectorielles}
% ============================================================

\subsection{Qu'est-ce qu'une image vectorielle ?}

\begin{quote}
\textbf{Une image vectorielle est une image composée de formes géométriques décrites mathématiquement : points, lignes, courbes. Chaque forme est définie par des coordonnées, des équations, des propriétés (couleur, contour).}
\end{quote}

Une image vectorielle n'est pas composée de pixels, mais de \terme{formes géométriques}. Ces formes sont décrites par des équations mathématiques, ce qui permet de les agrandir ou de les réduire à l'infini sans perte de qualité.

\subsection{Les éléments d'une image vectorielle}

\subsubsection{Le point d'ancrage}

Un \terme{point d'ancrage} (ou \emph{anchor point}) est un point qui définit une forme. Il a des coordonnées (x, y).

\subsubsection{Le segment}

Un \terme{segment} est une ligne qui relie deux points d'ancrage. Il peut être droit ou courbe.

\subsubsection{La courbe de Bézier}

Une \terme{courbe de Bézier} est une courbe définie par des points d'ancrage et des poignées de direction. Elle permet de créer des courbes précises et fluides.

\subsubsection{Le tracé}

Un \terme{tracé} (ou \emph{path}) est une suite de segments connectés. Il peut être ouvert ou fermé.

\subsubsection{Le contour}

Le \terme{contour} (ou \emph{stroke}) est la ligne visible d'un tracé. Il a une couleur, une épaisseur, un style.

\subsubsection{Le remplissage}

Le \terme{remplissage} (ou \emph{fill}) est la couleur ou le motif qui remplit l'intérieur d'une forme fermée.

\subsection{Avantages et inconvénients des images vectorielles}

\subsubsection{Avantages}

\begin{itemize}
    \item \terme{Redimensionnables} : pas de perte de qualité à l'agrandissement.
    \item \terme{Légères} : fichiers souvent plus légers que les images matricielles.
    \item \terme{Précises} : formes nettes, contours précis.
    \item \terme{Modifiables} : chaque forme peut être modifiée indépendamment.
    \item \terme{Compatibles} : formats universels (SVG, PDF, EPS, AI).
    \item \terme{Idéales} : pour les logos, les icônes, les illustrations.
\end{itemize}

\subsubsection{Inconvénients}

\begin{itemize}
    \item \terme{Photographie} : ne conviennent pas aux photographies.
    \item \terme{Dégradés complexes} : moins adaptées aux dégradés complexes.
    \item \terme{Effets} : moins adaptées aux effets de texture, de flou.
    \item \terme{Apprentissage} : plus difficiles à maîtriser que les images matricielles.
\end{itemize}

\subsection{Comparaison avec les images matricielles}

\begin{table}[h!]
\centering
\begin{tabular}{@{}p{4cm}p{5cm}p{5cm}@{}}
\toprule
\textbf{Critère} & \textbf{Image vectorielle} & \textbf{Image matricielle} \\
\midrule
Composition & Formes géométriques & Grille de pixels \\
Qualité à l'agrandissement & Aucune perte & Perte (pixellisation) \\
Poids & Souvent léger & Souvent lourd \\
Retouche & Non adaptée & Adaptée \\
Création de logos & Adaptée & Non adaptée \\
Logiciel & Illustrator & Photoshop \\
Formats & AI, SVG, EPS, PDF & JPEG, PNG, TIFF, PSD \\
\bottomrule
\end{tabular}
\caption{Comparaison image vectorielle / image matricielle.}
\end{table}

\subsection{Quand utiliser une image vectorielle ?}

\begin{itemize}
    \item \terme{Logos} : pour une qualité parfaite à toutes les tailles.
    \item \terme{Illustrations} : pour des formes nettes, précises.
    \item \terme{Icônes} : pour les interfaces, la signalétique.
    \item \terme{Pictogrammes} : pour les infographies, les schémas.
    \item \terme{Typographie} : pour les compositions typographiques.
    \item \terme{Infographies} : pour les graphiques, les schémas.
    \item \terme{Mises en page} : pour les affiches, les dépliants.
    \item \terme{Web} : pour les SVG, les icônes.
\end{itemize}

\subsection*{Conclusion de la section}

Une \terme{image vectorielle} est composée de formes géométriques décrites mathématiquement : points d'ancrage, segments, courbes de Bézier, tracés, contours, remplissages. Elle offre de nombreux avantages (redimensionnable, légère, précise, modifiable) mais a aussi des limites (pas adaptée à la photographie). Elle se distingue des images matricielles, qui sont composées de pixels. On utilise les images vectorielles pour les logos, les illustrations, les icônes, les pictogrammes, la typographie, les infographies.

Maintenant que nous comprenons les images vectorielles, nous allons, dans la section suivante, approfondir la \terme{différence entre images matricielles et vectorielles}.

% ============================================================
\section{Différence entre images matricielles et vectorielles}
% ============================================================

\subsection{Le principe fondamental}

La différence fondamentale entre les images matricielles et les images vectorielles réside dans leur \terme{mode de composition} :

\begin{itemize}
    \item \terme{Image matricielle} : composée de \emph{pixels}, chaque pixel ayant une couleur.
    \item \terme{Image vectorielle} : composée de \emph{formes géométriques}, chaque forme ayant des propriétés (coordonnées, couleur, contour).
\end{itemize}

Cette différence a des conséquences majeures sur la qualité, le poids, la modification, l'usage.

\subsection{La qualité à l'agrandissement}

\subsubsection{Image matricielle}

Lorsqu'on agrandit une image matricielle, les pixels deviennent visibles. C'est la \terme{pixellisation} : les bords deviennent crénelés, les détails se perdent, l'image devient floue.

\subsubsection{Image vectorielle}

Lorsqu'on agrandit une image vectorielle, les formes sont recalculées mathématiquement. La qualité reste parfaite : les bords sont nets, les détails sont préservés.

\subsubsection{Exemple}

Une même image (par exemple, un logo) en 100 × 100 pixels :

\begin{itemize}
    \item En \terme{matriciel} : agrandie à 1000 × 1000 pixels, elle sera pixellisée.
    \item En \terme{vectoriel} : agrandie à 1000 × 1000 pixels, elle sera parfaite.
\end{itemize}

\subsection{Le poids}

\subsubsection{Image matricielle}

Le poids d'une image matricielle dépend de sa résolution, de ses dimensions, de son contenu. Une image haute résolution peut être très lourde (plusieurs Mo).

\subsubsection{Image vectorielle}

Le poids d'une image vectorielle dépend du nombre de formes, de leur complexité. Une image simple peut être très légère (quelques Ko).

\subsubsection{Conséquence}

Pour le web, les images vectorielles (SVG) sont souvent plus légères et plus rapides à charger.

\subsection{La modification}

\subsubsection{Image matricielle}

La modification d'une image matricielle consiste à modifier les pixels. On peut :

\begin{itemize}
    \item \terme{Retoucher} : corriger les couleurs, les imperfections.
    \item \terme{Peindre} : ajouter des pixels.
    \item \terme{Effacer} : supprimer des pixels.
    \item \terme{Filtrer} : appliquer des effets.
\end{itemize}

\subsubsection{Image vectorielle}

La modification d'une image vectorielle consiste à modifier les formes. On peut :

\begin{itemize}
    \item \terme{Déplacer} : déplacer les points d'ancrage.
    \item \terme{Déformer} : modifier les courbes.
    \item \terme{Changer} : changer la couleur, le contour, le remplissage.
    \item \terme{Combiner} : combiner plusieurs formes.
\end{itemize}

\subsection{Les usages}

\subsubsection{Images matricielles}

\begin{itemize}
    \item \terme{Photographie} : captation, retouche, composition.
    \item \terme{Photomontage} : assemblage de photos.
    \item \terme{Illustration numérique} : peinture numérique.
    \item \terme{Textures} : création de textures réalistes.
    \item \terme{Effets} : flou, lumière, ombre.
\end{itemize}

\subsubsection{Images vectorielles}

\begin{itemize}
    \item \terme{Logos} : qualité parfaite à toutes les tailles.
    \item \terme{Illustrations} : formes nettes, précises.
    \item \terme{Icônes} : interfaces, signalétique.
    \item \terme{Pictogrammes} : infographies, schémas.
    \item \terme{Typographie} : compositions, lettrage.
    \item \terme{Mises en page} : affiches, dépliants.
    \item \terme{Web} : SVG, icônes.
\end{itemize}

\subsection{La complémentarité}

\subsubsection{Le principe}

Les images matricielles et vectorielles ne s'opposent pas : elles sont \terme{complémentaires}. Un designer professionnel utilise les deux, en fonction du besoin.

\subsubsection{Exemples}

\begin{itemize}
    \item \terme{Logo} : créé en vectoriel (Illustrator), exporté en PNG pour le web.
    \item \terme{Affiche} : composée d'images matricielles (photos) et d'éléments vectoriels (logo, textes).
    \item \terme{Infographie} : composée de formes vectorielles et d'images matricielles.
    \item \terme{Packaging} : logo vectoriel, photos de produit matricielles.
\end{itemize}

\subsection{Les passerelles entre les deux}

\subsubsection{Du vectoriel au matriciel}

On peut \terme{convertir} une image vectorielle en image matricielle (export en JPEG, PNG, TIFF). C'est une opération courante pour le web ou l'impression.

\subsubsection{Du matriciel au vectoriel}

On peut \terme{vectoriser} une image matricielle (transformer les pixels en formes). C'est une opération plus complexe, qui donne des résultats variables :

\begin{itemize}
    \item \terme{Images simples} : la vectorisation peut être bonne (logos, icônes).
    \item \terme{Images complexes} : la vectorisation est souvent médiocre (photos).
\end{itemize}

\subsection{Les bonnes pratiques}

\begin{itemize}
    \item \terme{Choisir} le bon type d'image selon l'usage.
    \item \terme{Utiliser} le vectoriel pour les logos, les icônes, les illustrations.
    \item \terme{Utiliser} le matriciel pour les photos, les textures.
    \item \terme{Convertir} si nécessaire (vectoriel → matriciel).
    \item \terme{Vectoriser} avec précaution (matriciel → vectoriel).
    \item \terme{Combiner} les deux types dans un projet.
    \item \terme{Respecter} les formats d'export.
\end{itemize}

\subsection*{Conclusion de la section}

La \terme{différence entre images matricielles et vectorielles} réside dans leur mode de composition : pixels pour les matricielles, formes géométriques pour les vectorielles. Cette différence a des conséquences sur la qualité à l'agrandissement, le poids, la modification, l'usage. Les images vectorielles sont adaptées aux logos, aux icônes, aux illustrations ; les images matricielles aux photos, aux textures. Les deux types sont complémentaires et peuvent être combinés dans un projet.

Maintenant que nous comprenons la différence entre les deux types d'images, nous allons, dans la section suivante, découvrir l'\terme{interface d'Illustrator}.

% ============================================================
\section{Découverte de l'interface}
% ============================================================

\subsection{Pourquoi maîtriser l'interface ?}

L'interface d'Illustrator est le \terme{poste de travail} du designer. Elle regroupe tous les outils, tous les panneaux, tous les menus dont tu auras besoin. La maîtriser est indispensable pour travailler efficacement.

Au début, l'interface d'Illustrator peut paraître intimidante, mais elle est en réalité très logique, une fois qu'on en comprend l'organisation.

\subsection{Les grandes zones de l'interface}

L'interface d'Illustrator se compose de plusieurs zones principales.

\subsubsection{La barre de menus}

Située en haut de l'écran, la \terme{barre de menus} contient tous les menus du logiciel : Fichier, Édition, Objet, Texte, Sélection, Effet, Affichage, Fenêtre, Aide.

\subsubsection{La barre de contrôle}

Située juste sous la barre de menus, la \terme{barre de contrôle} affiche les options de l'outil actuellement sélectionné, ainsi que les propriétés de l'objet sélectionné.

\subsubsection{La barre d'outils}

Située à gauche de l'écran (par défaut), la \terme{barre d'outils} contient tous les outils d'Illustrator : sélection, plume, formes, pinceau, texte, etc.

\subsubsection{Les panneaux}

Situés à droite de l'écran (par défaut), les \terme{panneaux} regroupent les fonctionnalités par thème : Propriétés, Calques, Couleur, Nuancier, Contour, Caractère, etc.

\subsubsection{Le plan de travail}

Au centre, le \terme{plan de travail} (ou \emph{artboard}) est la zone de création. C'est là que l'on dessine, compose, met en page.

\subsubsection{La barre d'état}

En bas de la fenêtre, la \terme{barre d'état} affiche des informations sur le document : niveau de zoom, etc.

\subsection{Organisation par défaut}

Par défaut, l'interface d'Illustrator est organisée ainsi :

\begin{itemize}
    \item \terme{En haut} : barre de menus, barre de contrôle.
    \item \terme{À gauche} : barre d'outils.
    \item \terme{À droite} : panneaux (Propriétés, Calques, Couleur, etc.).
    \item \terme{Au centre} : plan de travail.
    \item \terme{En bas} : barre d'état.
\end{itemize}

Cette organisation peut être modifiée selon les besoins.

\subsection{Personnalisation de l'interface}

\subsubsection{Déplacer les panneaux}

Les panneaux peuvent être déplacés par \terme{glisser-déposer}. On peut les détacher, les attacher, les empiler, les regrouper.

\subsubsection{Réduire les panneaux}

On peut réduire les panneaux en icônes pour gagner de la place, en cliquant sur la double flèche en haut à droite du panneau.

\subsubsection{Espaces de travail}

Illustrator propose plusieurs espaces de travail prédéfinis (\texttt{Fenêtre > Espace de travail}) :

\begin{itemize}
    \item \terme{Essentiel} : espace par défaut.
    \item \terme{Automatisation} : pour les scripts.
    \item \terme{Disposition} : pour la mise en page.
    \item \terme{Peinture} : pour l'illustration.
    \item \terme{Présentation} : pour les présentations.
    \item \terme{Impression et épreuvage} : pour l'impression.
    \item \terme{Tracé} : pour le dessin vectoriel.
    \item \terme{Typographie} : pour le texte.
    \item \terme{Web} : pour le web.
\end{itemize}

On peut aussi créer son propre espace de travail et l'enregistrer.

\subsubsection{Thème de l'interface}

Illustrator propose deux thèmes : \terme{clair} et \terme{sombre}. On peut les changer dans \texttt{Édition > Préférences > Interface} (Windows) ou \texttt{Illustrator > Préférences > Interface} (macOS). Le thème sombre est souvent préféré pour le design, car il réduit la fatigue visuelle et évite les interférences avec la perception des couleurs.

\subsection{Le plan de travail}

\subsubsection{Définition}

Le \terme{plan de travail} (ou \emph{artboard}) est la zone de création. C'est là que l'on dessine, compose, met en page. Il correspond à la zone qui sera exportée ou imprimée.

\subsubsection{Les plans de travail multiples}

Illustrator permet de créer plusieurs plans de travail dans un même document. Cela permet de :

\begin{itemize}
    \item \terme{Créer} plusieurs versions d'un visuel.
    \item \terme{Préparer} plusieurs formats (affiche, flyer, etc.).
    \item \terme{Organiser} les éléments d'un projet.
\end{itemize}

\subsubsection{Gérer les plans de travail}

On peut :

\begin{itemize}
    \item \terme{Créer} un nouveau plan de travail (\texttt{Fenêtre > Plans de travail}).
    \item \terme{Modifier} un plan de travail (position, dimensions).
    \item \terme{Renommer} un plan de travail.
    \item \terme{Supprimer} un plan de travail.
    \item \terme{Réorganiser} les plans de travail.
\end{itemize}

\subsection{Les raccourcis d'interface}

\begin{itemize}
    \item \terme{Tab} : masquer/afficher tous les panneaux.
    \item \terme{Maj + Tab} : masquer/afficher les panneaux sauf la barre d'outils.
    \item \terme{F} : basculer entre les modes d'affichage (standard, plein écran avec menus, plein écran).
    \item \terme{Ctrl + 0} : ajuster le plan de travail à l'écran.
    \item \terme{Ctrl + 1} : afficher à taille réelle.
\end{itemize}

\subsection{Les bonnes pratiques}

\begin{itemize}
    \item \terme{Choisir} un thème adapté (sombre pour le design).
    \item \terme{Organiser} ses panneaux en fonction de son usage.
    \item \terme{Enregistrer} son espace de travail personnalisé.
    \item \terme{Utiliser} les raccourcis pour gagner du temps.
    \item \terme{Ne pas hésiter} à réinitialiser l'espace de travail si nécessaire.
    \item \terme{Maîtriser} les plans de travail multiples.
\end{itemize}

\subsection*{Conclusion de la section}

L'\terme{interface d'Illustrator} se compose de plusieurs zones : barre de menus, barre de contrôle, barre d'outils, panneaux, plan de travail, barre d'état. Elle peut être personnalisée : déplacement des panneaux, réduction en icônes, choix d'un espace de travail, choix du thème. Le plan de travail est la zone de création : on peut en créer plusieurs dans un même document. Maîtriser l'interface permet de travailler efficacement et de se concentrer sur la création.

Maintenant que nous avons découvert l'interface, nous allons, dans la section suivante, étudier les \terme{plans de travail}.

% ============================================================
\section{Les plans de travail}
% ============================================================

\subsection{Qu'est-ce qu'un plan de travail ?}

\begin{quote}
\textbf{Un plan de travail (ou \emph{artboard}) est une zone de création délimitée, qui correspond à la zone qui sera exportée ou imprimée. Un document Illustrator peut contenir plusieurs plans de travail.}
\end{quote}

Le plan de travail est l'équivalent de la « page » dans un logiciel de mise en page. C'est la zone dans laquelle on crée le visuel, et qui détermine les dimensions du fichier final.

\subsection{Pourquoi utiliser plusieurs plans de travail ?}

\subsubsection{Créer plusieurs versions}

On peut créer plusieurs plans de travail pour :

\begin{itemize}
    \item \terme{Présenter} plusieurs versions d'un logo.
    \item \terme{Proposer} plusieurs pistes créatives.
    \item \terme{Décliner} un visuel en plusieurs formats.
\end{itemize}

\subsubsection{Préparer plusieurs formats}

On peut créer plusieurs plans de travail pour :

\begin{itemize}
    \item \terme{Préparer} une affiche (A3), un flyer (A5), une carte de visite.
    \item \terme{Décliner} un visuel pour les réseaux sociaux (carré, portrait, story).
    \item \terme{Organiser} les éléments d'un projet.
\end{itemize}

\subsubsection{Organiser un projet}

On peut créer plusieurs plans de travail pour :

\begin{itemize}
    \item \terme{Présenter} les éléments d'une identité visuelle (logo, couleurs, typographies).
    \item \terme{Présenter} les déclinaisons (papeterie, signalétique, web).
    \item \terme{Présenter} la charte graphique.
\end{itemize}

\subsection{Le panneau Plans de travail}

\subsubsection{Accès au panneau}

Le panneau \terme{Plans de travail} est accessible par \texttt{Fenêtre > Plans de travail}.

\subsubsection{Les options du panneau}

Le panneau propose :

\begin{itemize}
    \item \terme{Nouveau plan de travail} : crée un nouveau plan.
    \item \terme{Dupliquer} : duplique le plan sélectionné.
    \item \terme{Supprimer} : supprime le plan sélectionné.
    \item \terme{Réorganiser} : réorganise les plans.
    \item \terme{Options} : accède aux options du plan (nom, dimensions, orientation).
\end{itemize}

\subsection{Créer un plan de travail}

\subsubsection{Par le panneau Plans de travail}

\begin{enumerate}
    \item Ouvrir le panneau \emph{Plans de travail}.
    \item Cliquer sur \emph{Nouveau plan de travail}.
    \item Définir les dimensions, l'orientation, le nom.
    \item Valider.
\end{enumerate}

\subsubsection{Par l'outil Plan de travail}

\begin{enumerate}
    \item Sélectionner l'outil \emph{Plan de travail} (dans le groupe de l'outil Déplacement).
    \item Dessiner un rectangle pour définir le plan.
    \item Ajuster les dimensions.
\end{enumerate}

\subsubsection{Par le menu}

\texttt{Objet > Plans de travail > Insérer un nouveau plan de travail}.

\subsection{Modifier un plan de travail}

\subsubsection{Modifier les dimensions}

\begin{enumerate}
    \item Sélectionner l'outil \emph{Plan de travail}.
    \item Cliquer sur le plan.
    \item Ajuster les dimensions dans la barre de contrôle, ou en tirant sur les poignées.
\end{enumerate}

\subsubsection{Modifier la position}

\begin{enumerate}
    \item Sélectionner l'outil \emph{Plan de travail}.
    \item Cliquer sur le plan.
    \item Glisser le plan pour le déplacer.
\end{enumerate}

\subsubsection{Renommer un plan}

\begin{enumerate}
    \item Ouvrir le panneau \emph{Plans de travail}.
    \item Double-cliquer sur le nom du plan.
    \item Saisir un nouveau nom.
    \item Valider.
\end{enumerate}

\subsection{Les options de plan de travail}

\subsubsection{Le nom}

Le \terme{nom} du plan est important : il apparaîtra lors de l'export, et permettra d'identifier chaque plan.

\subsubsection{Les dimensions}

Les \terme{dimensions} (largeur et hauteur) définissent la taille du plan. On peut choisir des dimensions prédéfinies (A4, A3, etc.) ou personnalisées.

\subsubsection{L'orientation}

L'\terme{orientation} peut être portrait ou paysage.

\subsubsection{Le fond}

Le \terme{fond} peut être blanc, transparent, ou d'une couleur.

\subsubsection{Les repères}

Les \terme{repères} (règles, grille, repères commentés) aident à positionner les éléments.

\subsection{L'export des plans de travail}

\subsubsection{Exporter tous les plans}

\texttt{Fichier > Exporter > Exporter sous}, puis choisir \emph{Tous les plans de travail}.

\subsubsection{Exporter un plan spécifique}

\texttt{Fichier > Exporter > Exporter sous}, puis choisir \emph{Plage} et sélectionner le plan.

\subsubsection{Les formats d'export}

\begin{itemize}
    \item \terme{PNG} : pour le web, avec transparence.
    \item \terme{JPEG} : pour le web, sans transparence.
    \item \terme{PDF} : pour l'impression.
    \item \terme{SVG} : pour le web, vectoriel.
\end{itemize}

\subsection{Les bonnes pratiques}

\begin{itemize}
    \item \terme{Nommer} chaque plan de travail.
    \item \terme{Organiser} les plans de manière logique.
    \item \terme{Utiliser} des plans de travail multiples pour les projets complexes.
    \item \terme{Définir} les bonnes dimensions dès le départ.
    \item \terme{Vérifier} l'export de chaque plan.
    \item \terme{Ne pas surcharger} un plan (un plan = un visuel).
    \item \terme{Enregistrer} régulièrement.
\end{itemize}

\subsection*{Conclusion de la section}

Le \terme{plan de travail} est la zone de création qui correspond au visuel final. Un document Illustrator peut contenir plusieurs plans de travail, ce qui permet de créer plusieurs versions, plusieurs formats, ou d'organiser un projet. On peut créer, modifier, renommer, organiser les plans de travail via le panneau Plans de travail ou l'outil Plan de travail. Il est essentiel de nommer les plans, d'organiser, de définir les bonnes dimensions, de vérifier l'export.

Maintenant que nous connaissons les plans de travail, nous allons, dans la section suivante, étudier les \terme{outils essentiels} d'Illustrator.

% ============================================================
\section{Les outils essentiels}
% ============================================================

\subsection{La barre d'outils}

La barre d'outils regroupe tous les outils d'Illustrator. Elle est organisée par catégories. Certains outils sont masqués derrière d'autres : on y accède en cliquant longuement sur l'icône, ou en utilisant le raccourci clavier correspondant.

\subsection{Les outils de sélection}

\subsubsection{L'outil Sélection (V)}

L'\terme{outil Sélection} (flèche noire) permet de sélectionner et de déplacer des objets. On peut :

\begin{itemize}
    \item \terme{Sélectionner} un objet en cliquant dessus.
    \item \terme{Déplacer} un objet en le glissant.
    \item \terme{Redimensionner} un objet en tirant sur les poignées.
    \item \terme{Pivoter} un objet en plaçant le curseur à l'extérieur des poignées.
\end{itemize}

\subsubsection{L'outil Sélection directe (A)}

L'\terme{outil Sélection directe} (flèche blanche) permet de sélectionner et de modifier les points d'ancrage et les segments d'un objet. On peut :

\begin{itemize}
    \item \terme{Sélectionner} un point d'ancrage.
    \item \terme{Déplacer} un point d'ancrage.
    \item \terme{Modifier} les poignées de direction.
    \item \terme{Sélectionner} un segment.
\end{itemize}

\subsubsection{L'outil Lasso (Q)}

L'\terme{outil Lasso} permet de sélectionner des objets en dessinant une forme à main levée.

\subsubsection{L'outil Sélection de groupe}

Permet de sélectionner un groupe d'objets d'un seul clic.

\subsubsection{L'outil Sélection à l'intérieur}

Permet de sélectionner un objet à l'intérieur d'un groupe, sans sélectionner le groupe entier.

\subsection{Les outils de dessin}

\subsubsection{L'outil Plume (P)}

L'\terme{outil Plume} est l'outil de dessin vectoriel par excellence. Il permet de créer des tracés précis (courbes de Bézier). On peut :

\begin{itemize}
    \item \terme{Cliquer} pour créer un point d'angle.
    \item \terme{Cliquer-glisser} pour créer un point lisse.
    \item \terme{Modifier} les poignées de direction.
    \item \terme{Fermer} un tracé.
\end{itemize}

\subsubsection{L'outil Courbure}

L'\terme{outil Courbure} permet de dessiner des courbes lisses en cliquant-glissant.

\subsubsection{L'outil Crayon (N)}

L'\terme{outil Crayon} permet de dessiner à main levée. Le tracé est lissé automatiquement.

\subsubsection{L'outil Pinceau (B)}

L'\terme{outil Pinceau} permet de dessiner avec des formes de pinceau (calligraphie, art, motif, soies).

\subsubsection{L'outil Tache de peinture}

L'\terme{outil Tache de peinture} permet de créer des formes organiques en dessinant.

\subsection{Les outils de forme}

\subsubsection{L'outil Rectangle (M)}

L'\terme{outil Rectangle} permet de créer des rectangles. On peut :

\begin{itemize}
    \item \terme{Cliquer-glisser} pour dessiner.
    \item \terme{Maintenir Maj} pour un carré.
    \item \terme{Maintenir Alt} pour dessiner depuis le centre.
\end{itemize}

\subsubsection{L'outil Rectangle arrondi}

Permet de créer des rectangles à coins arrondis.

\subsubsection{L'outil Ellipse (L)}

L'\terme{outil Ellipse} permet de créer des ellipses et des cercles.

\subsubsection{L'outil Polygone}

Permet de créer des polygones réguliers (pentagone, hexagone, etc.).

\subsubsection{L'outil Étoile}

Permet de créer des étoiles.

\subsubsection{L'outil Spirale}

Permet de créer des spirales.

\subsubsection[L'outil Trait (\textbackslash)]{L'outil Trait (\textbackslash)}

Permet de créer des traits.

\subsection{Les outils de transformation}

\subsubsection{L'outil Rotation (R)}

L'\terme{outil Rotation} permet de faire pivoter des objets.

\subsubsection{L'outil Mise à l'échelle (S)}

L'\terme{outil Mise à l'échelle} permet de redimensionner des objets.

\subsubsection{L'outil Miroir (O)}

L'\terme{outil Miroir} permet de retourner des objets horizontalement ou verticalement.

\subsubsection{L'outil Transformation manuelle (E)}

L'\terme{outil Transformation manuelle} permet de transformer des objets librement.

\subsubsection{L'outil Largeur}

L'\terme{outil Largeur} permet de modifier l'épaisseur d'un tracé.

\subsubsection{L'outil Déformation (Maj + R)}

L'\terme{outil Déformation} permet de déformer des objets.

\subsubsection{L'outil Concepteur de forme (Maj + M)}

L'\terme{outil Concepteur de forme} permet de combiner, fusionner ou soustraire des formes.

\subsubsection{L'outil Ciseaux (C)}

L'\terme{outil Ciseaux} permet de couper un tracé.

\subsubsection{L'outil Couteau}

L'\terme{outil Couteau} permet de découper un objet en plusieurs parties.

\subsection{Les outils de peinture}

\subsubsection{L'outil Filet de dégradé (G)}

L'\terme{outil Filet de dégradé} permet de créer des dégradés en modifiant les mailles.

\subsubsection{L'outil Dégradé de formes}

L'\terme{outil Dégradé de formes} (Mélange) permet de créer une transition progressive entre deux formes.

\subsubsection{L'outil Pot de peinture dynamique}

L'\terme{outil Pot de peinture dynamique} permet de remplir des zones avec une couleur.

\subsubsection{L'outil Gomme (Maj + E)}

L'\terme{outil Gomme} permet d'effacer des parties d'objets.

\subsection{Les outils de texte}

\subsubsection{L'outil Texte (T)}

L'\terme{outil Texte} permet de créer du texte. On peut :

\begin{itemize}
    \item \terme{Cliquer} pour un texte ponctuel.
    \item \terme{Cliquer-glisser} pour un texte de paragraphe.
    \item \terme{Cliquer sur un tracé} pour un texte curviligne.
\end{itemize}

\subsubsection{L'outil Texte vertical}

Permet de créer du texte vertical.

\subsubsection{L'outil Texte sur tracé}

Permet de créer du texte sur un tracé.

\subsubsection{L'outil Texte curviligne}

Permet de créer du texte sur une courbe.

\subsection{Les outils de navigation}

\subsubsection{L'outil Main (H)}

L'\terme{outil Main} permet de déplacer la vue dans le plan de travail.

\subsubsection{L'outil Zoom (Z)}

L'\terme{outil Zoom} permet de zoomer avant ou arrière.

\subsection{Les outils de couleur}

\subsubsection{Le sélecteur de couleur}

Le \terme{sélecteur de couleur} permet de choisir une couleur.

\subsubsection{Le nuancier}

Le \terme{nuancier} permet de choisir une couleur dans une bibliothèque.

\subsection{Les raccourcis des outils}

\begin{table}[h!]
\centering
\begin{tabular}{@{}ll@{}}
\toprule
\textbf{Outil} & \textbf{Raccourci} \\
\midrule
Sélection & V \\
Sélection directe & A \\
Lasso & Q \\
Plume & P \\
Crayon & N \\
Pinceau & B \\
Rectangle & M \\
Ellipse & L \\
Rotation & R \\
Mise à l'échelle & S \\
Miroir & O \\
Transformation manuelle & E \\
Concepteur de forme & Maj + M \\
Ciseaux & C \\
Dégradé & G \\
Gomme & Maj + E \\
Texte & T \\
Main & H \\
Zoom & Z \\
\bottomrule
\end{tabular}
\caption{Raccourcis clavier des outils essentiels.}
\end{table}

\subsection{Les bonnes pratiques}

\begin{itemize}
    \item \terme{Apprendre} les raccourcis progressivement.
    \item \terme{Utiliser} l'outil adapté à la tâche.
    \item \terme{Ne pas hésiter} à expérimenter.
    \item \terme{Combiner} les outils pour obtenir le résultat souhaité.
    \item \terme{Travailler} de manière non destructive.
\end{itemize}

\subsection*{Conclusion de la section}

Illustrator propose de nombreux outils, classés par catégories : sélection, dessin, forme, transformation, peinture, texte, navigation, couleur. Chaque outil a une fonction spécifique et un raccourci clavier. Maîtriser les outils essentiels est indispensable pour travailler efficacement. Les raccourcis clavier permettent de gagner du temps et de se concentrer sur la création.

Maintenant que nous connaissons les outils essentiels, nous allons, dans la section suivante, étudier la \terme{création et l'enregistrement d'un document}.

% ============================================================
\section{Création et enregistrement d'un document}
% ============================================================

\subsection{Créer un nouveau document}

\subsubsection{La commande Nouveau}

Pour créer un nouveau document : \texttt{Fichier > Nouveau} (raccourci \texttt{Ctrl + N} sur Windows, \texttt{Cmd + N} sur macOS).

\subsubsection{La boîte de dialogue}

La boîte de dialogue \emph{Nouveau document} propose plusieurs paramètres.

\paragraph{Nom} : donner un nom au document.

\paragraph{Profil} : choisir un profil prédéfini (Impression, Web, Mobile, Film et vidéo, etc.).

\paragraph{Taille} : choisir une taille prédéfinie (A4, A3, etc.) ou personnalisée.

\paragraph{Largeur et hauteur} : définir les dimensions du plan de travail.

\paragraph{Unités} : pixels, centimètres, millimètres, pouces, points, picas.

\paragraph{Orientation} : portrait ou paysage.

\paragraph{Plans de travail} : nombre de plans de travail.

\paragraph{Espacement} : espace entre les plans de travail.

\paragraph{Mode colorimétrique} : RVB, CMJN.

\paragraph{Effet raster} : résolution pour les effets matriciels (72 ppi, 150 ppi, 300 ppi, etc.).

\paragraph{Aperçu} : afficher un aperçu.

\subsubsection{Les profils prédéfinis}

Illustrator propose de nombreux profils prédéfinis :

\begin{itemize}
    \item \terme{Impression} : A4, A3, A5, etc.
    \item \terme{Web} : 1920 × 1080, 1280 × 720, etc.
    \item \terme{Mobile} : iPhone, iPad, etc.
    \item \terme{Film et vidéo} : 4K, Full HD, etc.
    \item \terme{Art et illustration} : formats courants.
\end{itemize}

\subsubsection{Les bonnes pratiques}

\begin{itemize}
    \item \terme{Définir} les bonnes dimensions dès le départ.
    \item \terme{Choisir} le bon mode colorimétrique (RVB pour l'écran, CMJN pour l'impression).
    \item \terme{Choisir} la bonne résolution d'effet raster (300 ppi pour l'impression, 72 ppi pour le web).
    \item \terme{Nommer} le document de manière claire.
    \item \terme{Définir} le nombre de plans de travail si nécessaire.
\end{itemize}

\subsection{Ouvrir un document}

\subsubsection{La commande Ouvrir}

Pour ouvrir un document : \texttt{Fichier > Ouvrir} (raccourci \texttt{Ctrl + O} sur Windows, \texttt{Cmd + O} sur macOS).

\subsubsection{Les formats supportés}

Illustrator supporte de nombreux formats : AI, SVG, EPS, PDF, mais aussi des formats d'autres logiciels (PSD, JPEG, PNG, TIFF, etc.).

\subsubsection{L'ouverture récente}

Pour ouvrir un document récent : \texttt{Fichier > Ouvrir récent}.

\subsubsection{L'importation}

Pour importer un fichier dans un document existant : \texttt{Fichier > Importer} ou \texttt{Fichier > Importer et placer}.

\subsection{Enregistrer un document}

\subsubsection{La commande Enregistrer}

Pour enregistrer un document : \texttt{Fichier > Enregistrer} (raccourci \texttt{Ctrl + S} sur Windows, \texttt{Cmd + S} sur macOS).

Si le document est nouveau, la boîte de dialogue \emph{Enregistrer sous} s'ouvre.

\subsubsection{La commande Enregistrer sous}

Pour enregistrer sous un autre nom ou format : \texttt{Fichier > Enregistrer sous} (raccourci \texttt{Ctrl + Maj + S} sur Windows, \texttt{Cmd + Maj + S} sur macOS).

\subsubsection{Les formats d'enregistrement}

\begin{itemize}
    \item \terme{AI} : format natif d'Illustrator, préserve les calques.
    \item \terme{SVG} : format vectoriel pour le web.
    \item \terme{EPS} : format pour l'impression.
    \item \terme{PDF} : format de document.
\end{itemize}

\subsubsection{La commande Exporter}

Pour exporter un document : \texttt{Fichier > Exporter > Exporter sous} ou \texttt{Fichier > Exporter > Exporter pour les écrans}.

\paragraph{Exporter sous} : permet de choisir le format (PNG, JPEG, SVG, etc.), la qualité, les dimensions.

\paragraph{Exporter pour les écrans} : permet d'exporter les plans de travail en plusieurs formats et tailles.

\subsubsection{Les options d'enregistrement}

\begin{itemize}
    \item \terme{Créer un fichier PDF compatible} : permet d'ouvrir le fichier AI dans d'autres logiciels.
    \item \terme{Inclure les plans de travail} : enregistre tous les plans de travail.
    \item \terme{Inclure les polices} : incorpore les polices utilisées.
    \item \terme{Utiliser la compression} : réduit le poids du fichier.
\end{itemize}

\subsection{Les bonnes pratiques}

\begin{itemize}
    \item \terme{Toujours conserver le AI} comme fichier source.
    \item \terme{Enregistrer régulièrement} pour éviter les pertes.
    \item \terme{Utiliser la commande Enregistrer sous} pour créer des versions.
    \item \terme{Exporter dans le bon format} selon l'usage.
    \item \terme{Incorporer les polices} dans le PDF.
    \item \terme{Organiser ses fichiers} dans des dossiers clairs.
    \item \terme{Faire des sauvegardes} régulières.
\end{itemize}

\subsection*{Conclusion de la section}

Créer, ouvrir et enregistrer un document sont des opérations fondamentales dans Illustrator. La création passe par la boîte de dialogue \emph{Nouveau}, qui permet de définir les dimensions, le mode colorimétrique, la résolution d'effet raster, le nombre de plans de travail. L'ouverture se fait par la commande \emph{Ouvrir} ou \emph{Importer}. L'enregistrement se fait par les commandes \emph{Enregistrer}, \emph{Enregistrer sous}, \emph{Exporter}. Il est essentiel de connaître les formats et de choisir le bon selon l'usage.

Maintenant que nous savons créer, ouvrir et enregistrer un document, nous allons, dans la section suivante, étudier les \terme{raccourcis clavier fondamentaux}.

% ============================================================
\section{Les raccourcis clavier fondamentaux}
% ============================================================

\subsection{Pourquoi apprendre les raccourcis ?}

Les raccourcis clavier sont essentiels pour travailler efficacement dans Illustrator. Ils permettent de :

\begin{itemize}
    \item \terme{Gagner du temps} : éviter de chercher dans les menus.
    \item \terme{Rester concentré} : ne pas interrompre le travail.
    \item \terme{Améliorer la fluidité} : enchaîner les actions rapidement.
    \item \terme{Travailler comme un professionnel} : les pros utilisent tous les raccourcis.
\end{itemize}

Au début, il peut être difficile de les mémoriser. Mais avec la pratique, ils deviennent automatiques.

\subsection{Les raccourcis essentiels}

\subsubsection{Fichier}

\begin{table}[h!]
\centering
\begin{tabular}{@{}ll@{}}
\toprule
\textbf{Action} & \textbf{Raccourci} \\
\midrule
Nouveau document & Ctrl + N (Cmd + N) \\
Ouvrir & Ctrl + O (Cmd + O) \\
Enregistrer & Ctrl + S (Cmd + S) \\
Enregistrer sous & Ctrl + Maj + S (Cmd + Maj + S) \\
Exporter sous & Ctrl + Alt + Maj + S (Cmd + Opt + Maj + S) \\
Fermer & Ctrl + W (Cmd + W) \\
Imprimer & Ctrl + P (Cmd + P) \\
\bottomrule
\end{tabular}
\caption{Raccourcis de gestion de fichiers.}
\end{table}

\subsubsection{Édition}

\begin{table}[h!]
\centering
\begin{tabular}{@{}ll@{}}
\toprule
\textbf{Action} & \textbf{Raccourci} \\
\midrule
Annuler & Ctrl + Z (Cmd + Z) \\
Rétablir & Ctrl + Maj + Z (Cmd + Maj + Z) \\
Copier & Ctrl + C (Cmd + C) \\
Coller & Ctrl + V (Cmd + V) \\
Couper & Ctrl + X (Cmd + X) \\
Tout sélectionner & Ctrl + A (Cmd + A) \\
Désélectionner & Ctrl + Maj + A (Cmd + Maj + A) \\
\bottomrule
\end{tabular}
\caption{Raccourcis d'édition.}
\end{table}

\subsubsection{Outils}

\begin{table}[h!]
\centering
\begin{tabular}{@{}ll@{}}
\toprule
\textbf{Outil} & \textbf{Raccourci} \\
\midrule
Sélection & V \\
Sélection directe & A \\
Lasso & Q \\
Plume & P \\
Crayon & N \\
Pinceau & B \\
Rectangle & M \\
Ellipse & L \\
Rotation & R \\
Mise à l'échelle & S \\
Miroir & O \\
Transformation manuelle & E \\
Concepteur de forme & Maj + M \\
Ciseaux & C \\
Dégradé & G \\
Gomme & Maj + E \\
Texte & T \\
Main & H \\
Zoom & Z \\
\bottomrule
\end{tabular}
\caption{Raccourcis des outils.}
\end{table}

\subsubsection{Objets}

\begin{table}[h!]
\centering
\begin{tabular}{@{}ll@{}}
\toprule
\textbf{Action} & \textbf{Raccourci} \\
\midrule
Grouper & Ctrl + G (Cmd + G) \\
Dégrouper & Ctrl + Maj + G (Cmd + Maj + G) \\
Verrouiller & Ctrl + 2 \\
Déverrouiller tout & Ctrl + Alt + 2 \\
Masquer & Ctrl + 3 \\
Afficher tout & Ctrl + Alt + 3 \\
Placer devant & Ctrl + Maj + ] \\
Placer derrière & Ctrl + Maj + [ \\
Avancer & Ctrl + ] \\
Reculer & Ctrl + [ \\
\bottomrule
\end{tabular}
\caption{Raccourcis de gestion des objets.}
\end{table}

\subsubsection{Affichage}

\begin{table}[h!]
\centering
\begin{tabular}{@{}ll@{}}
\toprule
\textbf{Action} & \textbf{Raccourci} \\
\midrule
Zoom avant & Ctrl + + \\
Zoom arrière & Ctrl + - \\
Ajuster à l'écran & Ctrl + 0 (zéro) \\
Taille réelle & Ctrl + 1 \\
Masquer les panneaux & Tab \\
Masquer sauf la barre d'outils & Maj + Tab \\
Afficher/masquer les règles & Ctrl + R \\
Afficher/masquer les repères & Ctrl + ; \\
Afficher/masquer la grille & Ctrl + ' \\
\bottomrule
\end{tabular}
\caption{Raccourcis d'affichage.}
\end{table}

\subsubsection{Texte}

\begin{table}[h!]
\centering
\begin{tabular}{@{}ll@{}}
\toprule
\textbf{Action} & \textbf{Raccourci} \\
\midrule
Texte & T \\
Texte vertical & Maj + T \\
Texte sur tracé & Maj + T (sur un tracé) \\
Créer des contours & Ctrl + Maj + O (Cmd + Maj + O) \\
\bottomrule
\end{tabular}
\caption{Raccourcis de texte.}
\end{table}

\subsection{Les bonnes pratiques}

\begin{itemize}
    \item \terme{Apprendre} progressivement : ne pas essayer de tout mémoriser d'un coup.
    \item \terme{Pratiquer} régulièrement : utiliser les raccourcis dès que possible.
    \item \terme{Créer} ses propres raccourcis : Illustrator permet de personnaliser les raccourcis (\texttt{Édition > Raccourcis clavier}).
    \item \terme{Afficher} les raccourcis : on peut afficher les raccourcis dans les menus.
    \item \terme{Ne pas hésiter} à utiliser les menus au début.
\end{itemize}

\subsection*{Conclusion de la section}

Les \terme{raccourcis clavier} sont essentiels pour travailler efficacement dans Illustrator. Ils couvrent la gestion de fichiers, l'édition, les outils, les objets, l'affichage, le texte. Il faut les apprendre progressivement et les pratiquer régulièrement. Illustrator permet également de personnaliser ses raccourcis.

Maintenant que nous connaissons les raccourcis fondamentaux, nous allons, dans la dernière section de ce chapitre, étudier les \terme{formats AI, SVG, EPS et PDF}.

% ============================================================
\section{Les formats AI, SVG, EPS et PDF}
% ============================================================

\subsection{Pourquoi connaître les formats ?}

Connaître les formats de fichiers est essentiel pour :

\begin{itemize}
    \item \terme{Choisir} le format adapté à l'usage (web, impression, édition).
    \item \terme{Préserver} la qualité du fichier.
    \item \terme{Optimiser} le poids du fichier.
    \item \terme{Éviter} les incompatibilités.
    \item \terme{Préserver} la possibilité de modifier le fichier.
\end{itemize}

\subsection{Le format AI}

\subsubsection{Définition}

Le format \terme{AI} (Adobe Illustrator) est le format natif d'Illustrator. Il préserve tous les calques, les plans de travail, les tracés, les effets, les polices.

\subsubsection{Caractéristiques}

\begin{itemize}
    \item \terme{Format natif} d'Illustrator.
    \item \terme{Préserve} tous les éléments (calques, plans de travail, tracés, effets).
    \item \terme{Modifiable} : on peut rouvrir le fichier et continuer à travailler.
    \item \terme{Poids} variable (dépend du contenu).
    \item \terme{Compatible} uniquement avec Illustrator (et quelques autres logiciels Adobe).
\end{itemize}

\subsubsection{Usages}

\begin{itemize}
    \item \terme{Fichiers de travail} Illustrator.
    \item \terme{Archivage} de projets.
    \item \terme{Échanges} entre designers.
\end{itemize}

\subsubsection{Précautions}

\begin{itemize}
    \item \terme{Toujours conserver} le AI comme fichier source.
    \item \terme{Ne pas livrer} le AI au client (sauf demande explicite).
    \item \terme{Exporter} en SVG, PDF, PNG pour la diffusion.
\end{itemize}

\subsection{Le format SVG}

\subsubsection{Définition}

Le format \terme{SVG} (Scalable Vector Graphics) est un format vectoriel basé sur XML, standard du web. Il est utilisé pour les logos, les icônes, les illustrations vectorielles.

\subsubsection{Caractéristiques}

\begin{itemize}
    \item \terme{Vectoriel} : pas de perte de qualité à l'agrandissement.
    \item \terme{Léger} : fichiers souvent plus légers que les images matricielles.
    \item \terme{Modifiable} : on peut modifier le code XML.
    \item \terme{Compatible} avec les navigateurs modernes.
    \item \terme{Animable} : on peut animer les SVG avec CSS ou JavaScript.
    \item \terme{Accessible} : les SVG peuvent être lus par les lecteurs d'écran.
\end{itemize}

\subsubsection{Usages}

\begin{itemize}
    \item \terme{Logos} : pour les sites web.
    \item \terme{Icônes} : pour les interfaces.
    \item \terme{Illustrations} : pour les pages web.
    \item \terme{Infographies} : pour les articles.
    \item \terme{Cartes} : interactives.
\end{itemize}

\subsubsection{Précautions}

\begin{itemize}
    \item \terme{Éviter} les effets complexes (flou, ombre) : mal supportés.
    \item \terme{Vectoriser} les textes : pour éviter les problèmes de polices.
    \item \terme{Optimiser} le fichier : supprimer les éléments inutiles.
\end{itemize}

\subsection{Le format EPS}

\subsubsection{Définition}

Le format \terme{EPS} (Encapsulated PostScript) est un format vectoriel basé sur le langage PostScript. Il a longtemps été le standard pour l'impression et l'échange de fichiers.

\subsubsection{Caractéristiques}

\begin{itemize}
    \item \terme{Vectoriel} : pas de perte de qualité.
    \item \terme{Universel} : compatible avec de nombreux logiciels.
    \item \terme{PostScript} : adapté à l'impression professionnelle.
    \item \terme{Aperçu} : contient une image matricielle d'aperçu.
    \item \terme{Poids} variable.
\end{itemize}

\subsubsection{Usages}

\begin{itemize}
    \item \terme{Impression} professionnelle.
    \item \terme{Échanges} entre professionnels.
    \item \terme{Archivage} de fichiers vectoriels.
\end{itemize}

\subsubsection{Précautions}

\begin{itemize}
    \item \terme{Moins moderne} que le PDF : préférer le PDF pour l'impression.
    \item \terme{Moins compatible} que le SVG pour le web.
    \item \terme{Utiliser} l'EPS si nécessaire (anciens logiciels, imprimeurs spécifiques).
\end{itemize}

\subsection{Le format PDF}

\subsubsection{Définition}

Le format \terme{PDF} (Portable Document Format) est un format de document développé par Adobe. Il peut contenir du texte, des images vectorielles, des images matricielles, des liens, des formulaires.

\subsubsection{Caractéristiques}

\begin{itemize}
    \item \terme{Multi-plateforme} : lisible sur tous les appareils.
    \item \terme{Préserve} la mise en page.
    \item \terme{Peut contenir} du texte, des images, des liens.
    \item \terme{Peut être imprimé} fidèlement.
    \item \terme{Peut être protégé} par mot de passe.
\end{itemize}

\subsubsection{Les variantes du PDF}

\begin{itemize}
    \item \terme{PDF/X-1a} : standard pour l'impression, CMJN, polices incorporées.
    \item \terme{PDF/X-4} : standard moderne, transparence gérée.
    \item \terme{PDF/A} : standard pour l'archivage.
    \item \terme{PDF interactif} : avec formulaires, liens, animations.
\end{itemize}

\subsubsection{Usages}

\begin{itemize}
    \item \terme{Documents} professionnels.
    \item \terme{Livraison} de fichiers à l'imprimeur.
    \item \terme{Présentations}.
    \item \terme{Portfolios}.
\end{itemize}

\subsection{Tableau récapitulatif des formats}

\begin{table}[h!]
\centering
\begin{tabular}{@{}p{2cm}p{3cm}p{2cm}p{4.5cm}@{}}
\toprule
\textbf{Format} & \textbf{Type} & \textbf{Modifiable} & \textbf{Usage principal} \\
\midrule
AI & Vectoriel & Oui & Fichiers de travail Illustrator \\
SVG & Vectoriel & Oui & Logos, icônes, web \\
EPS & Vectoriel & Oui & Impression, échanges \\
PDF & Mixte & Oui & Documents, impression \\
\bottomrule
\end{tabular}
\caption{Récapitulatif des principaux formats d'Illustrator.}
\end{table}

\subsection{Comment choisir un format ?}

\subsubsection{Pour le web}

\begin{itemize}
    \item \terme{Logos, icônes} : SVG.
    \item \terme{Illustrations} : SVG.
    \item \terme{Images matricielles} : PNG, JPEG (si nécessaire).
\end{itemize}

\subsubsection{Pour l'impression}

\begin{itemize}
    \item \terme{Documents} : PDF/X-1a ou PDF/X-4.
    \item \terme{Échanges} : EPS ou PDF.
    \item \terme{Fichiers de travail} : AI.
\end{itemize}

\subsubsection{Pour l'archivage}

\begin{itemize}
    \item \terme{Fichiers de travail} : AI.
    \item \terme{Documents} : PDF/A.
\end{itemize}

\subsection{Les bonnes pratiques}

\begin{itemize}
    \item \terme{Conserver} le fichier AI comme source.
    \item \terme{Choisir} le bon format selon l'usage.
    \item \terme{Vectoriser} les textes pour le SVG et l'EPS.
    \item \terme{Incorporer} les polices dans le PDF.
    \item \terme{Optimiser} le poids des SVG.
    \item \terme{Vérifier} le rendu du fichier exporté.
    \item \terme{Organiser} les fichiers exportés.
\end{itemize}

\subsection*{Conclusion de la section}

Illustrator propose plusieurs formats d'enregistrement et d'exportation : AI (format natif), SVG (format web), EPS (format d'échange), PDF (format de document). Chacun a ses caractéristiques et ses usages. Il est essentiel de conserver le AI comme source, de choisir le bon format selon l'usage, de vectoriser les textes pour le SVG et l'EPS, d'incorporer les polices dans le PDF, d'optimiser le poids des SVG.

% ============================================================
\section*{Conclusion du chapitre}
\addcontentsline{toc}{section}{Conclusion du chapitre}
% ============================================================

Nous voici au terme du chapitre 14. Nous avons découvert Adobe Illustrator sous tous ses aspects :

\begin{itemize}
    \item \terme{Présentation} : logiciel de création d'images vectorielles, standard professionnel.
    \item \terme{Domaines d'utilisation} : logos, identité visuelle, illustration, pictogrammes, typographie, infographies, mise en page, web, motion design, packaging.
    \item \terme{Images vectorielles} : composées de formes géométriques, redimensionnables sans perte.
    \item \terme{Différence matriciel/vectoriel} : pixels vs formes, qualité à l'agrandissement, poids, usages.
    \item \terme{Interface} : barre de menus, barre de contrôle, barre d'outils, panneaux, plan de travail.
    \item \terme{Plans de travail} : zones de création, plans multiples, gestion, export.
    \item \terme{Outils essentiels} : sélection, dessin, forme, transformation, peinture, texte, navigation.
    \item \terme{Création et enregistrement} : nouveau document, ouverture, enregistrement, export.
    \item \terme{Raccourcis clavier} : gain de temps, fluidité, professionnalisme.
    \item \terme{Formats} : AI, SVG, EPS, PDF ; caractéristiques et usages.
\end{itemize}

Illustrator est un outil puissant, mais il n'est qu'un outil. La maîtrise d'Illustrator ne remplace ni la culture visuelle, ni la créativité, ni la connaissance des principes du design. Elle les prolonge, elle les amplifie. Un designer qui maîtrise Illustrator peut créer des logos, des illustrations, des identités visuelles avec précision et qualité. Un designer qui maîtrise Illustrator sans culture visuelle ne produira que des formes vides.

Ce chapitre t'a donné les bases nécessaires pour utiliser Illustrator. Mais la maîtrise vient avec la pratique. N'hésite pas à expérimenter, à explorer, à te tromper, à recommencer. Chaque projet est une occasion d'apprendre.

Dans le chapitre suivant, nous aborderons le \terme{dessin vectoriel et les formes géométriques}. Tu apprendras à comprendre les tracés vectoriels, à utiliser les outils Rectangle et Ellipse, à créer des polygones et des étoiles, à modifier les dimensions des formes, à utiliser les outils de transformation, à aligner et distribuer les objets, à dupliquer et organiser les éléments, à comprendre les contours et les fonds, à combiner des formes géométriques, à créer des compositions simples. Ces compétences te permettront de dessiner avec précision dans Illustrator.

\vspace{2em}

\begin{center}
\textit{Fin du chapitre 14 --- Introduction à Adobe Illustrator.}

\vspace{1em}

\textbf{Prochain chapitre : Chapitre 15 --- Dessin vectoriel et formes géométriques.}
\end{center}